\documentclass[
	a4paper,
	10pt,
	unnumberedsections,
	twoside, 
]{LTJournalArticle}

\documentclass[11pt,a4paper]{article}


\addbibresource{Sample.bib}

\runninghead{Communication in Distributed Architecture}

\footertext{\textit{SWE4913 Research Report} (2026)}

\setcounter{page}{1}

\usepackage{listings}
\usepackage{xcolor}
\usepackage{tabularx}
\usepackage{graphicx}
\usepackage[table]{xcolor}
\usepackage{amsmath}

\lstset{
	language=C++,
	basicstyle=\ttfamily\small,
	breaklines=true,
	frame=single,
	captionpos=b,
	keywordstyle=\color{blue},
	commentstyle=\color{green!50!black},
	stringstyle=\color{orange},
	numbers=left,
	numberstyle=\tiny\color{gray}	
}



%----------------------------------------------------------------------------------------
%	TITLE SECTION
%----------------------------------------------------------------------------------------

\title{SWE4913 \\Communication in Distributed Architecture\\  Optimizing Point-to-Point Data Transfer}
\author{%
	Eddie Ogendo
}

% Affiliations are output in the \date{} command
\date{\footnotesize Software Engineering, University of New Brunswick.}

% Full-width abstract
\renewcommand{\maketitlehookd}{%
	\begin{abstract}
		\noindent This research investigates the upper theoretical and practical limits of point-to-point data transfer throughput. While high-level frameworks like the Message Passing Interface (MPI) provide robust abstractions for distributed systems, they often introduce overhead through multiple memory-copy operations and context switching. We evaluate the impact of communication protocols, message sizes, and Modular Component Architecture (MCA) tuning, and propose a high-performance "Broadcaster" utilizing Linux io\_uring and zero-copy splice/tee primitives. Our results demonstrate that bypassing the user-space buffer bottleneck allows data transfer rates to approach the physical saturation point of the underlying hardware interconnect.
	\end{abstract}
}

%----------------------------------------------------------------------------------------

\begin{document}
\maketitle

\section{Introduction}
\label{sec:intro}

The movement of high-volume data remains a primary bottleneck in modern distributed simulation environments as the underlying communication architecture must handle gigabytes of mesh and field data with minimal latency. 

In the context of the \textit{Stage2} execution pipeline, operations are often coupled in a producer-consumer relationship. Traditional data transfer methods rely on high-level abstractions like the Message Passing Interface (\texttt{MPI}) or standard \texttt{POSIX} I/O. While robust, these methods introduce a "copy tax", the redundant migration of data between kernel and user-space memory, which exhausts the memory bus and consumes CPU cycles that would otherwise be dedicated to physics computation.

This research investigates the theoretical upper limits of data transfer by bypassing these user-space bottlenecks. We propose a "Broadcaster" architecture built on the Linux \texttt{io\_uring} subsystem. By utilizing kernel-level primitives such as \texttt{splice()} and \texttt{tee()}, we demonstrate a zero-copy mechanism that duplicates data streams directly within kernel pipe-buffers. This approach effectively shifts the performance bottleneck from software-defined overhead to the physical saturation limits of the host hardware.



% --- Section 2: High Performance Parallelism ---
\section{High Performance Parallelism}
\label{sec:hpp}
\subsection{Message Passing Interface (MPI)}
\label{sec:MPI}
The Message Passing Interface (\texttt{MPI}) serves as the standard for communication in distributed-memory systems. Its primary contribution to high-performance computing (\texttt{HPC}) is the abstraction of hardware heterogeneity and network complexity.

\begin{itemize}
	\item \textbf{Abstraction of Connectivity:} \texttt{MPI} use communicators and groups to isolate process communication. This allows complex simulations, such as \textit{FVCOM} ocean models to run across thousands of nodes without manually managing the socket connection
	
	\item \textbf{Layered Architecture:} In implementations like \texttt{OpenMPI}, the Modular Component Architecture (MCA) allows for runtime selection of transport layers. The Point-to-Point Management Layer (PML) handles message matching and ordering, while the Byte Transport Layer (BTL) interfaces with the physical hardware (e.g., InfiniBand via UCX or Ethernet via TCP).
\end{itemize}

\subsection{The "Memory Wall" and Copy Overhead}
Despite its robustness, \texttt{MPI} introduces a significant "copy tax." In an \texttt{MPI\_Send(...)} operation, data often undergoes multiple transformations:

\begin{enumerate}
	\item \textbf{Application Buffer to MPI Buffer:} Data is copied into an internal buffer for "eager" or "rendezvous" protocol handling.
	\item \textbf{User-space to Kernel-space:} The data is passed to the OS network stack.
	\item \textbf{Kernel to Hardware:} The NIC finally pulls the data from the kernel buffer.
\end{enumerate}

This triple-copy approach is a manifestation of the \textbf{Memory Wall}, where the CPU's ability to process data outstrips the memory bus's ability to move it. In the context of \textit{Stage2}, where operations may reside on the same physical host, this overhead becomes the primary bottleneck.

\subsection{Open Multi-Processing (OpenMP)}
\label{sec:OpenMP}
An application programming interface (API) that supports multi-platform shared-memory multiprocessing programming in C, C++, and Fortran. 
It allows developers to write parallel code easily and efficiently by adding simple compiler directives to their existing code.

\subsection{Shared Memory vs. Distributed Memory}
\label{sec:OMPvsMPI}
While MPI handles distributed systems, OpenMP is typically used for shared-memory parallelism.
\begin{itemize}
\item \textbf{Thread-based Parallelism:} OpenMP operates on a single memory space, avoiding message passing entirely.
\item \textbf{Hybrid Constraints:} While OpenMP is faster for local tasks, it cannot scale across nodes. This creates a performance gap: MPI is too heavy for local "Operation" coupling, and OpenMP is too limited for network-wide distribution.
\end{itemize}
\textit{Differences between OpenMP and MPI: Table\ref{tab:mpi_vs_openmp}} 


\begin{table*}[ht] % 't' helps placement in twoside documents
	\centering
	\small % Slightly smaller text helps with dense comparison tables
	\begin{tabularx}{\textwidth}{|X|X|} % 'X' columns auto-size to fill the space
		\hline
		\rowcolor{gray!10} % Optional: requires \usepackage[table]{xcolor}
		\textbf{MPI (Message Passing Interface)} & \textbf{OpenMP (Open Multi-Processing)} \\ \hline
		Supports parallel computation for distributed and shared-memory systems. & Supports parallel computation for shared-memory systems only. \\ \hline
		Process-based parallelism; processes are independent. & Thread-based parallelism; threads share resources. \\ \hline
		Each process has its own private memory space. & Threads share the same address space and memory. \\ \hline
		Data is exchanged explicitly via message passing. & Data is shared implicitly; no explicit message passing. \\ \hline
		Initialization overhead usually occurs once at startup. & Thread creation overhead can occur dynamically during tasks. \\ \hline
	\end{tabularx}
	\caption{Architectural Comparison: MPI vs. OpenMP}
	\label{tab:mpi_vs_openmp}
\end{table*}

\subsection{Task / Channel Models in \textit{Stage2}}
Task / Channel Model represents computations as a set of tasks that may interact with each other by sending messages through channels. 

\textbf{Tasks} can send local data values to other tasks via output ports and receive data from other tasks via input ports. Receiving in the model is \textbf{synchronous} while sending is a \textbf{asynchronous} operation.

\textbf{Channel} is a message queue that connects one task's output with another task's input port. Data values appear at the input port in the same order as which they were placed in the output port at the end of the channel.

Foster's design encourages development of scalable parallel algorithms by delaying machine-dependent considerations to the later steps.The design's fours steps involves:
\begin{enumerate}
	\item Partitioning.
	\item Communication.
	\item Agglomeration.
	\item Mapping
\end{enumerate}

We can incorporate Foster’s design methodology by treating \textit{Stage2} operations as discrete tasks and the kernel-level pipes as channels allowing reimplementing the "Collective Communication" patterns (Broadcast, Scatter, Gather) using zero-copy primitives rather than high-level buffered messages.


\section{Networking Foundations}
\label{sec:Network}

\subsection{Resumable / Partial Downloads}
Partial or resumable downloads allow a client to retrieve only specific portions of a resource rather than downloading the entire file from the beginning. This is essential for large media files, software packages, and any scenario where network interruptions are common.  
To support partial downloads, the \textbf{server / sender} must explicitly implement HTTP range semantics and advertise that it can return byte ranges on demand.

\subsection{Resumable / Multipart Uploads}
Resumable or multipart uploads are designed to reliably transfer large files over unstable or bandwidth‑limited networks. Instead of sending the entire file in a single request, the file is divided into smaller, manageable chunks that are uploaded independently. This approach improves reliability, reduces memory usage, and allows uploads to continue even after interruptions.

\textbf{Resumability} is achieved by tracking which chunks have been successfully acknowledged by the receiver.  
If a failure occurs; whether due to network instability, client crash, or server timeout, the sender retransmits only the missing chunks rather than restarting the entire upload.

\subsection{Implementation of Resumability and Reliability}
Drawing from the semantics of HTTP Range Requests, our broadcaster treats the source file as a series of addressable byte-ranges.

\begin{itemize}
	\item \textbf{Partial Content (206 Logic):} By tracking offsets and the file descriptors within the \textit{io\_uring} SQEs, we ensure that if a specific transfer chunk fails, only that "range" needs to be re-dispatched to the specific file.
	\item \textbf{Error Isolation:} Similar to Multipart Uploads, a failure in one broadcast "sub-pipe" (one destination) does not cancel the entire operation for other nodes depending on the error handling protocol implemented.
\end{itemize}

\subsection{The Logic of Chunking}
\label{sec:chunking}
To saturate a network link between two Stage2 nodes, data cannot be sent as a single monolithic block. We utilize a chunk-based transfer strategy (specifically between $64\text{ KB}$ to $1\text{ MB}$ in the implementation). This size is chosen to align with the default Linux pipe capacity, minimizing "pipe stalls" while remaining large enough to amortize the overhead of the io\_uring submission.


\subsection{Latency}
\label{sec:latency}
Time it takes for data to go from the origin to the destination. This delay encompasses various factors such as:
\begin{itemize}
	\item \textbf{Processing delay} - time it takes for data to be processed by network devices like routers \& switches
	\item \textbf{Queueing delay} - time spent in queue waiting to be processed by the to the destination host.
	\item \textbf{Transmission delays} - time it takes to push the packet’s bits into the wire.
	\item \textbf{Propagation delays} - time the signal takes to propagate through the medium. It would be a signal through an optic cable.
\end{itemize}
It refers to network’s responsiveness and perceived speed and it is measured in milliseconds. The lower the "latency time" in milliseconds the better

\begin{flushleft}
	\begin{align*}
		&\text{Packet size } (L) = \text{Packet size} \\
		&\text{Average rate of arrival } (a) = \frac{\text{packets}}{\text{sec}} \\
		&\text{Bandwidth } (R) = \frac{b}{\text{sec}} \\
		&\text{Traffic intensity} = L \cdot \frac{a}{R} \\
		&\text{Propagation delay } (d_{\text{prop}}) = \frac{m}{s} \\
		&\text{Transmission delay } (d_{\text{trans}}) = \frac{L}{R} \\
		&\rho = \frac{1}{R} \\
		&\text{Time of 1st bit arrival (ms)} = d_{\text{prop}} + \rho \\
		&\text{Total Delay} = d_{\text{proc}} + d_{\text{queue}} + d_{\text{trans}} + d_{\text{prop}} \\
	\end{align*}
\end{flushleft}


\label{sec:background}
\subsection{Jitter}
The measure of the variation of latency over time, measured in \textit{milliseconds}. It is the measure of the variability in packet arrival times, which can be caused by:
\begin{itemize}
	\item \textbf{Network congestion}: High traffic volumes causing packets to take longer paths or experience delays.
	\item \textbf{Routing Changes}: Dynamic changes in the path that data packets take across the network.
	\item \textbf{Variable Path Lengths}: Differences in the physical distance data packets must travel.
	\item \textbf{Equipment Limitations}: Older routers with limited processing power struggle to handle modern network traffic volumes. When a router's buffer fills up, it must decide which packets to prioritize and which to delay or drop entirely.
\end{itemize}
It affects the quality of real-time communication. If the jitter is high then it may cause packet loss, poor audio and video quality transmission, and data streaming interruptions.Thus, jitter is a significant objective to consider in the performance of networks serving real-time applications.


\begin{figure*} 
	\includegraphics[width=\linewidth]{speed_test.jpeg}
	\caption{Internet speed test on IPv4(connection) from Fredericton(Client location) to Newark(Server location). Source cloudflare, \href{https://speed.cloudflare.com/}{https://speed.cloudflare.com/}.}
	\label{fig:InternetSpeedTest}
\end{figure*}

\begin{table*}[ht] % 't' helps placement in twoside documents
	\centering
	\small % Slightly smaller text helps with dense comparison tables
	\begin{tabularx}{\textwidth}{|X|X|X|} % 'X' columns auto-size to fill the space
		\hline
		\rowcolor{gray!10} % Optional: requires \usepackage[table]{xcolor}
		\textbf{Aspect} & \textbf{Latency} & \textbf{Jitter} \\	\hline
		Definition & Time taken for data to travel from source to destination & Variation in latency over time \\ \hline
		Measurement & Typically measured in milliseconds (ms) & Also measured in milliseconds (ms) \\ \hline
		Components & Includes processing delay, queuing delay, transmission delay, propagation delay, etc. & Reflects inconsistency in latency due to network fluctuations \\	\hline
		Impact & Affects overall speed and responsiveness of network connections & Impacts consistency and stability of data transmission \\	\hline
		Nature & Generally a static value & Dynamic and can fluctuate over time \\ \hline
		Examples & 50 ms latency means it takes 50 ms for data to travel from source to destination & Jitter of 5 ms means latency varies by up to 5 ms over time \\ \hline
	\end{tabularx}
	\caption{Comparison between Latency and Jitter}
	\label{tab:latency_vs_jitter}
\end{table*}


In distributed architecture, the time to first bit arrival is governed by \textbf{Total Delay}. The goal is to minimize $d_{\text{proc}}$ (processing delay) by using Zero-Copy techniques, eliminating the CPU time spent on moving bits from a network buffer to an application buffer, we reduce the overall \textit{Jitter} of the simulation pipeline thus, bringing the observed bandwidth closer to the theoretical maximum $R$.

\section{Kernel-level Optimization}
\subsection{Concurrency Architecture}
Applications dealing with millions of requests in a small time frame and are concerned with a efficiency, use asynchronous programming to avoid the operating system’s thread/process context switching overhead by processing much of the I/O in a single thread. 

\subsubsection{Iterative:}
This processes one request after another. While it is serves one request, other requests that might arrive have to wait till the previous one is done processing. 
Some system call would tend to have a limit, like \href{http://man7.org/linux/man-pages/man2/listen.2.html}{\texttt{listen(2)}}, to how many requests the operating system will queue up.

\subsubsection{Forking:}
This creates a new process for each request that needs to be served. Requests don’t have to wait for previous requests to get processed. 

\subsubsection{Preforked:}
Avoids the overhead of having to create a totally new process every time a request needs to be served by initializing a pool of processes that are assigned requests as they come in. Only when all the pooled processes are busy that, incoming requests have to wait for their turn to get processed. The number of processes can be adjusted depending on the load usually experienced.

\subsubsection{Threaded:}
This spawns a new thread every time a request needs to be processed. Threads share a lot of data with the main process creating it and thus incur a slightly lower overhead during creation compared to creating a new process.

\subsubsection{Prethreaded:}
This is the threads equivalent of the `pre-forked` architecture. In this style, a pool of threads are created and from the pool requests are assigned as they come in. As with the pre-forked model, requests have to wait only if all threads are busy processing previously received requests.It is very efficient model commonly followed by most web application frameworks.

\subsubsection{Poll:}
A single threaded and uses the \href{http://man7.org/linux/man-pages/man2/poll.2.html}{\texttt{poll(2)}} system call to multiplex between requests. However, \href{http://man7.org/linux/man-pages/man2/poll.2.html}{\texttt{poll(2)}} has performance problems when scaling to a large number of file descriptors. The state for each request is tracked and a series of callbacks to functions are made that take processing of that request to the next stage.

\subsection{IO\_Uring}
\label{sec:io_uring}
Linux-specific API for asynchronous I/O. It allows the user to submit one or more I/O requests, which are processed asynchronously without blocking the calling process.
Removes the limitations of the current family of system calls used to monitor file descriptors(e.g. sockets) and let you know when one or more of them are ready.
\begin{itemize}
	\item \href{http://man7.org/linux/man-pages/man2/select.2.html}{\texttt{select(2)}} allows a program to monitor multiple file descriptors but, can only monitor file descriptors numbers that are less than \texttt{FD\_SETSIZE}(1024)
	\item  \href{http://man7.org/linux/man-pages/man2/select.2.html}{\texttt{select(2)}}, \href{http://man7.org/linux/man-pages/man2/poll.2.html}{\texttt{poll(2)}} or the \href{http://man7.org/linux/man-pages/man7/select.7.html}{\texttt{epoll()}} family of system calls will always return regular files as being ready for I/O. Linux kernel treats disk files as \textit{always ready} even though the actual data might not be in the CPU's RAM (Page Cache) but in the ssd or hdd.
\end{itemize}

\texttt{io\_uring} moves away from "Readiness" (Is it ready?) to "Completion" (Tell me when you're done). It instructs the kernel to get the data from the disk and put it in this memory buffer and only alert it if the data is actually there.

\subsection{Asynchronous Orchestration via io\_uring}
To reach the hardware's throughput limits, the software should reduce the overhead of synchronous system calls. Traditional I/O requires a context switch for every operation, stalling the CPU during disk or network wait times. This implementation of the Stage2 Broadcaster uses \texttt{io\_uring}, which utilizes a shared-memory ring buffer architecture. This allows the application to place multiple I/O requests in a Submission Queue (SQ) and move on to other tasks, only interacting with the kernel again to reap results from the Completion Queue (CQ).

For \textit{Stage2} operations, such as distributing an \textit{OpenFOAM} mesh, this means the broadcaster can prepare a full 1MB chunk transfer to $N$ consumers in a single kernel entry, significantly reducing the Processing Delay ($d_{proc}$) defined in \ref{sec:latency}{latency}.

\subsection{Zero-Copy Stream Duplication: Splice and Tee}
The primary driver of this research is the elimination of the "User-Space Copy Tax." By utilizing \texttt{splice()} and \texttt{tee()} primitives, data is never mapped into the application's address space. 
\begin{itemize}
	\item \textbf{\texttt{splice()}}: Facilitates the movement of data between a file descriptor and a kernel pipe-buffer by manipulating page-table references rather than copying bits.
	\item \textbf{\texttt{tee()}}: A non-consumptive operation that "clones" the contents of the master pipe into multiple sub-pipes without depleting the source buffer.
\end{itemize}
This architectural choice allows the broadcaster to achieve high-efficiency "Data Fan-out," where the CPU acts as an orchestrator of kernel memory pointers rather than a mover of data, preserving cycles for the computational operations.

\subsection{Managing Pipe Capacity and Deadlocks}
A critical challenge identified during testing was the "Pipe Deadlock." Because \texttt{tee()} duplicates data without consuming it, the master pipe remains filled after the duplication phase, halting subsequent read operations. 

To resolve this, we implemented a \textbf{Terminal Consumption Splice}. The final destination in the broadcast group ($N^{th}$ node) does not utilize \texttt{tee()}; instead, it performs a standard \texttt{splice()} from the master pipe. This "clears the drain," freeing kernel pages for the next $64\text{ KB}$ chunk. To further optimize throughput, \texttt{fcntl()} was utilized with \texttt{F\_SETPIPE\_SZ} to expand pipe buffers to $1\text{ MB}$. This alignment with our chunking strategy maximizes the "flight time" of data, ensuring the hardware remains saturated between \texttt{io\_uring} submissions.

\subsection{Solving the Linked-Chain Fragility}
A second technical challenge involved managing the \texttt{io\_uring} engine's tendency for out-of-order execution. Without strict sequencing, the kernel might attempt to consume data via a destination \texttt{splice} before the \texttt{tee} duplication is complete, leading to data corruption or \texttt{EAGAIN} errors.

Initially, a monolithic asynchronous batch submission was attempted using \texttt{IOSQE\_IO\_LINK}. However, this resulted in widespread \texttt{ECANCELED} errors; if any single destination in the chain experienced a short-write or stall, the kernel would discard the entire subsequent broadcast for all other nodes. 

We addressed this by implementing a \textbf{Hybrid-Link Submission Model}. In this model:
\begin{enumerate}
	\item The initial \texttt{source} $\rightarrow$ \texttt{master-pipe} read is submitted as an independent, unlinked SQE.
	\item For the first $N-1$ destinations, we submit linked pairs of (\texttt{tee} $\rightarrow$ \texttt{splice}). This ensures the destination write only triggers upon a successful duplication.
	\item The $N^{th}$ destination is submitted as a standalone \texttt{splice} to drain the master pipe. 
\end{enumerate}
This model effectively isolates failure domains; a bottleneck in one consumer no longer interrupts the data flow to the remaining nodes in the Stage2 execution pipeline.


\subsection{Implementation of the Broadcaster Loop}
The resulting logic follows a strict batch-and-wait cycle, ensuring that hardware saturation is reached before the next submission is issued.
\begin{lstlisting}[caption=Optimized Broadcaster Loop Logic]

io_uring_sqe *sqe = io_uring_get_sqe(&ring);
io_uring_prep_splice(/*...*/);
sqe->flags |= IOSQE_IO_LINK;
submissions++;
	
//  BROADCAST TO N-1 FILES
for (size_t i = 0; i < output_files.size() - 1; i++)
{
	// TEE: Source Pipe -> Dest Pipe (No offsets)
	sqe = io_uring_get_sqe(&ring);
	io_uring_prep_tee(/*...*/);
	sqe->flags |= IOSQE_IO_LINK;
	submissions++;
		
	// SPLICE: Dest Pipe -> File (Must use file offset)
	sqe = io_uring_get_sqe(&ring);
	io_uring_prep_splice(/*...*/);
	sqe->flags |= IOSQE_IO_LINK;
	submissions++;
}
	
sqe = io_uring_get_sqe(&ring);
io_uring_prep_splice(/*...*/);
submissions++;
	
total_processed += to_process;
\end{lstlisting}


\section{Synthesis and Benchmarking}
\subsection{Experimental Configuration}
To evaluate the efficacy of the zero-copy broadcaster, benchmarks were conducted within a virtualized Linux environment (\texttt{WSL2}) on a host equipped with high-speed NVMe storage. This setup specifically tests the Memory Bus and Virtual I/O throughput, simulating the high-speed backplane communication required when multiple \textit{Stage2} operations reside on the same physical node.

\subsection{Throughput vs. Chunk Size}
The relationship between chunk size and total bandwidth was analyzed to determine the optimal "flight size" for data in the kernel pipes.
\begin{itemize}
	\item \textbf{Small Chunks (4 KB - 32 KB):} Observed throughput was limited by the frequency of \texttt{io\_uring} completions and the overhead of managing a high volume of SQEs.
	\item \textbf{Optimal Chunks (64 KB - 1 MB):} By expanding the pipe capacity via \texttt{F\_SETPIPE\_SZ} to 1 MB, we achieved a significant increase in bandwidth. This alignment minimizes the "stalling" of the producer, allowing the $d_{trans}$ (transmission delay) to dominate over the $d_{proc}$ (processing delay).
\end{itemize}

\subsection{Comparative Analysis: Zero-Copy vs. Buffered I/O}
The zero-copy broadcaster was compared against a traditional buffered implementation (simulating standard MPI-style communication).

\begin{itemize}
\item \textbf{CPU Utilization:} While the buffered approach reached high throughput, it exhibited linear increases in CPU load due to the "Copy Tax."
\item \textbf{Saturation Point:} The \texttt{io\_uring} broadcaster reached the hardware's saturation point with significantly lower CPU overhead, proving that the bottleneck was shifted successfully from the software stack to the physical I/O limits.
\end{itemize}

\subsection{Application within Stage2 Execution Pipeline}
The benchmarking results confirm that this communication design meets the requirements for Stage2's "geometry-level" connection. For a typical simulation mesh (multi-gigabyte range), the broadcaster can distribute data to $N$ concurrent operations with near-zero latency added per additional consumer. This "Data Fan-out" capability is essential for real-time co-simulation where solvers, loggers, and analysis tools must synchronize without degrading the solver's performance.Analysis of the Benchmarking Logic


\section{Performance Evaluation and Discussion}
In this section, we analyze the performance of four distinct broadcasting strategies implemented using io\_uring. The evaluation focuses on throughput and execution time when replicating a 512 MB source file across multiple destinations.

\subsection{Experimental Setup}
The benchmarks were conducted on a Linux system (Kernel 6.x) using a 512 MB source file. Each mode was executed for four trials, with the system page cache cleared between runs via drop\_caches to ensure results reflect raw I/O performance rather than memory caching. The four tested modes are defined as follows:
\begin{itemize}
\item \textbf{Mode 1:} Standard readv/writev using OpenMP for parallel distribution.
\item \textbf{Mode 2:} io\_uring with vector I/O (readv/writev) and manual buffer management.
\item \textbf{Mode 3:} io\_uring utilizing splice for zero-copy data movement via intermediate pipes.
\item \textbf{Mode 4:} io\_uring utilizing tee and splice for zero-copy broadcasting.
\end{itemize}

\subsection{Comparative Analysis}
The data in table \ref{tab:comp_analysis} summarizes the average throughput and completion times across all trials.

\begin{table*}[ht] % 't' helps placement in twoside documents
	\centering
	\small % Slightly smaller text helps with dense comparison tables
	\begin{tabularx}{\textwidth}{|X|X|X|X|} % 'X' columns auto-size to fill the space
		\hline
		\rowcolor{gray!10} % Optional: requires \usepackage[table]{xcolor}
		\textbf{Strategy} & \textbf{Avg. Throughput (MB/s)} & \textbf{Avg. Time (s)} & \textbf{Efficiency Rank} \\ \hline
		OMP preadv/pwritev & 432.12 & 0.0624 & 4 \\ \hline 
		IO\_uring readv/writev & 439.29 & 0.0387 & 3 \\ \hline
		IO\_uring Splice & 461.91 & 0.0329 & 2 \\ \hline
		IO\_uring Tee + Splice & 611.25 & 0.0273 & 1 \\ \hline
	\end{tabularx}
	\caption{Comparison between different broadcasters}
	\label{tab:comp_analysis}
\end{table*}

\subsection{Analysis of User-Space Buffering (Modes 1 \& 2)}

\textbf{Mode 1} showed the highest variance, with throughput starting from \textit{137 MB/s} before settling near \textit{450 MB/s}. This fluctuations are likely due to the overhead of synchronization and the cost of \texttt{posix\_memalign} during the allocation phase in \texttt{readv\_broadcaster.hpp}. 

\textbf{Mode 2} (IO\_uring readv/writev) demonstrated a significant improvement in stability and a \textit{31\%} increase in average throughput over \textbf{Mode 1}, validating the efficiency of \texttt{io\_uring's} submission queue mechanism for handling multiple I/O vectors without the overhead of threading.

\subsection{The Zero-Copy Advantage (Modes 3 \& 4)}
\textbf{Modes 3 and 4} outperformed all buffer-based methods. Mode 3 which utilizes the \texttt{splice()} system call achieved an average of \textit{461.91 MB/s}. By moving data between two file descriptors without copying between kernel address space and user address space, this mode avoids the CPU cycles required to copy data onto user-space memory.

\textbf{Mode 4} (Tee + Splice) emerged as the most efficient strategy, achieving a peak throughput of \textit{620.62 MB/s}. This represents a \textit{\textit{~31\%}} improvement over Mode 3 and a \textit{~81\%} improvement over the baseline (Mode 1). The performance leap is attributed to the \texttt{tee()} system call, which allows the kernel to duplicate data within pipe buffers. This enables a single "read" from the source to feed multiple "writes" to destinations simultaneously, effectively bypassing the memory bandwidth bottleneck associated with traditional data replication.

\subsection{Technical Challenges and Optimization}
Two critical optimizations were used during development to stabilize these results:
\begin{enumerate}
\item \textbf{Pipe Buffer Saturation:} In \texttt{broadcastTee.hpp}, a CHUNK\_SIZE of 256 KB was utilized. Initial tests with larger chunks resulted in "Operation Canceled" errors due to pipe capacity limits.

\item \textbf{Linked SQEs:} The implementation utilized IOSQE\_IO\_LINK to chain the \texttt{tee} and \texttt{splice} operations. This ensures that the data duplication and the final "drain" from the source pipe occur in a strict, atomic-like sequence, preventing data loss or race conditions within the kernel pipes. This also allowed parallelization of tee requests as the linked splice ensure the pipe is empty on the next \texttt{tee}.
\end{enumerate}

\subsection{Summary}
The data confirms that for broadcasting tasks, minimizing the need for context switching from data between kernel and user space, and elimination of redundant data duplication are the most effective optimization. The io\_uring framework, when paired with tee and splice, provides a robust, high-performance pipeline that significantly exceeds the capabilities of standard multi-threaded I/O.

%----------------------------------------------------------------------------------------
%	 REFERENCES
%----------------------------------------------------------------------------------------

\begin{thebibliography}{9}
	\bibitem{openmpi} 
	The Open MPI Project. \textit{Open MPI: A High Performance Message Passing Library}. \textit{https://www.open-mpi.org/}
	
	\bibitem{https://www.youtube.com/watch?v=c0C9mQaxsD4} MPI Basics by Tom Nurkkalla. Introduction to distributed computing with MPI. \textit{https://www.youtube.com/}
	
	\bibitem{https://arcca.github.io/intro-mpi/} Introduction to Parallel Programming using MPI. \textit{https://arcca.github.io/intro-mpi/}
	
	\bibitem{https://rookiehpc.org/} RookieHPC, learn about High-Performance Computing. \textit{https://rookiehpc.org/}
	
	\bibitem{https://www.geeksforgeeks.org/computer-networks/latency-vs-jitter-in-computer-networks/} Latency vs Jitter in Computer Networks. \textit{https://www.geeksforgeeks.org/}
	
	\bibitem{https://speed.cloudflare.com/}Cloudflare internet speed test. \textit{https://speed.cloudflare.com/}
	
	\bibitem{https://en.wikipedia.org/wiki/CAP_theorem}CAP theorem \textit{https://en.wikipedia.org/wiki/CAP\_theorem}
	
	\bibitem{https://umamahesh.net/batch-vs-stream-processing-a-comprehensive-comparison-of-paradigms-with-use-cases/}Batch vs Stream Processing: \textit{A comprehensive comparison of paradigm}: \textit{https://umamahesh.net/batch-vs-stream-processing-a-comprehensive-comparison-of-paradigms-with-use-cases/}
	
	\bibitem{https://www.datacamp.com/blog/batch-vs-stream-processing}Batch vs Stream Processing: \textit{When to Use Each and Why It Matters}. \textit{https://www.datacamp.com/blog/batch-vs-stream-processing}
	
	\bibitem{https://www.amazon.com/Parallel-Programming-C-MPI-OpenMP/dp/0070582017?keywords=Parallel+Programming+in+C+With+Mpi+and+Openmp&qid=1684601986&s=books&sr=1-1&linkCode=sl1&tag=freecomputerbooks-20&linkId=5272329a690999da59b8b67c5a0705ab&language=en_US}Parallel programming in C with MPI and openMP \textit{by Michael J. Quinn}
	
\end{thebibliography}

\printbibliography % Output the bibliography

%----------------------------------------------------------------------------------------

\end{document}
